\chapter{測地線と変位補間}
\label{ch:seminar-07-geodesics}

第~\ref{ch:seminar-05-wasserstein}章で，最適輸送が測度の空間に距離を定めることを見た．
本章からは，その距離空間の\textbf{幾何}——とくに\textbf{測地線}——を扱う．
測地線は「質量を中間点へ連続的に動かす」操作なので，
地の空間を $\R^d$ のような連続空間にとる必要がある
（有限ビンの上では測地線が存在しない）．
そこで本章では連続設定を最小限に導入し，
最適輸送が定める測地線——\textbf{McCann の変位補間}——を構成する．
途中で必要となる「最適カップリングは凸関数の勾配 $\nabla\phi$ で与えられる」
という事実（Brenier）は，証明せず補題として引用する
（その役割は後の曲率の章で本質化する）．

%% ============================================================
%% 7.1 連続設定と 2-Wasserstein 距離
%% ============================================================
\section{連続設定と 2-Wasserstein 距離}
\label{sec:sem-continuous-w2}

地の空間を $\R^d$，地のコストを $c(x,y) = \norm{x-y}^2$（2 次コスト）とする．
2 次モーメントが有限な確率測度の空間
\[
  \Pp_2(\R^d)
  \defeq
  \left\{
    \alpha \in \Mm_+^1(\R^d)
    \;\middle|\;
    \int_{\R^d} \norm{x}^2 \,\d\alpha(x) < \infty
  \right\}
\]
の上で議論する．

\begin{definition}{2-Wasserstein 距離}{sem-w2-def}
  $\alpha, \beta \in \Pp_2(\R^d)$ に対して，\textbf{2-Wasserstein 距離}を
  \[
    \Wass_2(\alpha, \beta)
    \defeq
    \left(
      \inf_{\pi \in \Couplings(\alpha, \beta)}
      \int_{\R^d \times \R^d} \norm{x - y}^2 \,\d\pi(x, y)
    \right)^{1/2}
  \]
  で定める（連続 Kantorovich 問題；定義~\ref{def:sem-kantorovich}）．
\end{definition}

\begin{remark}{$\Wass_2$ は距離である}{sem-w2-is-metric}
  $\Wass_2$ は $\Pp_2(\R^d)$ 上の距離である．
  対称性・非退化性，および貼り合わせ補題による三角不等式の証明は，
  離散の場合（定理~\ref{thm:sem-wasserstein-is-distance}）と同じ論法で，
  和を積分に，明示構成
  $\mathbf{S} = \mathbf{P}\diag(1/\tilde{\mathbf{b}})\mathbf{Q}$ を
  分解（disintegration）による貼り合わせに置き換えればよい．
  また下半連続なコストに対し，$\Pp_2(\R^d)$ 上で最適カップリングが
  存在する（弱位相でのコンパクト性による）．
\end{remark}

%% ============================================================
%% 7.2 最適写像（Brenier の事実）
%% ============================================================
\section{最適写像：Brenier の事実}
\label{sec:sem-brenier-fact}

測地線を写像で書くために，最適カップリングが\textbf{写像}で与えられるための
事実を引用する．これは第~\ref{ch:seminar-06-dual}章の双対・$c$-変換・相補性の
連続版の結実であり，証明は与えない．

\begin{claim}{Brenier の定理}{sem-brenier}
  $\alpha \in \Pp_2(\R^d)$ が Lebesgue 測度に関して絶対連続
  （密度を持つ，$\alpha \ll \Lcal^d$）とし，$\beta \in \Pp_2(\R^d)$ とする．
  このとき $\Wass_2(\alpha,\beta)$ を達成する最適カップリングはただ一つで，
  ある\textbf{凸関数} $\phi : \R^d \to \R$ の勾配
  $T = \nabla\phi$ による\textbf{写像}で与えられる：
  \[
    \pi = (\Id, \nabla\phi)_\sharp \alpha,
    \qquad\text{すなわち}\qquad
    \beta = (\nabla\phi)_\sharp \alpha.
  \]
  この $T = \nabla\phi$ を $\alpha$ から $\beta$ への\textbf{最適輸送写像}とよぶ．
\end{claim}

\begin{remark}{ch06 との対応}{sem-brenier-from-dual}
  Brenier 写像が凸関数の勾配となるのは，
  第~\ref{ch:seminar-06-dual}章の構造の連続版である．
  2 次コストでは双対ポテンシャルの $c$-変換が\textbf{Legendre 変換}に対応し，
  相補性条件（命題~\ref{prop:sem-complementary-slackness}）
  $\supp(\pi) \subseteq \{f(x)+g(y)=c(x,y)\}$ が，
  $\phi(x) = \tfrac{\norm{x}^2}{2} - f(x)$ とおくと
  $y \in \partial\phi(x)$（劣微分）と書ける．
  $\alpha$ が密度を持てば $\phi$ は $\alpha$-ほとんど至るところ微分可能で，
  $y = \nabla\phi(x)$ が一意に定まる．これが写像 $T = \nabla\phi$ である．
  「凸関数の勾配」は，1 次元の単調増加関数の多次元への一般化にあたる．
\end{remark}

%% ============================================================
%% 7.3 変位補間と測地線
%% ============================================================
\section{変位補間と測地線}
\label{sec:sem-displacement-interpolation}

距離空間の測地線をまず定義する．

\begin{definition}{定速測地線}{sem-constant-speed-geodesic}
  距離空間 $(M, d_M)$ において，曲線
  $(\mu_t)_{t \in [0,1]}$ が端点 $\mu_0, \mu_1$ を結ぶ\textbf{定速測地線}であるとは，
  任意の $s, t \in [0,1]$ に対して
  \[
    d_M(\mu_s, \mu_t) = |t - s|\, d_M(\mu_0, \mu_1)
  \]
  が成り立つことをいう．
\end{definition}

Brenier 写像 $T = \nabla\phi$ に沿って質量を一定速度で動かすと，測地線が得られる．

\begin{definition}{McCann の変位補間}{sem-mccann-interpolation}
  $\alpha \ll \Lcal^d$，$\beta \in \Pp_2(\R^d)$ とし，
  $T = \nabla\phi$ を $\alpha$ から $\beta$ への最適輸送写像
  （主張~\ref{clm:sem-brenier}）とする．
  $t \in [0,1]$ に対し $T_t \defeq (1-t)\Id + t\,T$ とおき，
  \[
    \mu_t \defeq (T_t)_\sharp \alpha
    = \bigl((1-t)\Id + t\,\nabla\phi\bigr)_\sharp \alpha
  \]
  を\textbf{変位補間}とよぶ．$\mu_0 = \alpha$，$\mu_1 = \beta$ である．
\end{definition}

\begin{theorem}{変位補間は測地線である}{sem-mccann-geodesic}
  変位補間 $(\mu_t)_{t \in [0,1]}$ は，$(\Pp_2(\R^d), \Wass_2)$ における
  $\alpha$ と $\beta$ を結ぶ定速測地線である：
  \[
    \Wass_2(\mu_s, \mu_t) = |t - s|\, \Wass_2(\alpha, \beta)
    \qquad (\forall\, s, t \in [0,1]).
  \]
\end{theorem}

\begin{proof}
  まず最適写像の定義より，$\nabla\phi$ が $\alpha$ から $\beta$ への
  最適カップリングを与えるので
  \[
    \Wass_2(\alpha, \beta)^2
    = \int_{\R^d} \norm{x - \nabla\phi(x)}^2 \,\d\alpha(x).
    \tag{$\ast$}
  \]

  \textbf{上界．}
  $s \leq t$ とする．$(T_s, T_t)_\sharp \alpha$ は
  $\mu_s = (T_s)_\sharp\alpha$ と $\mu_t = (T_t)_\sharp\alpha$ の
  カップリングであるから，その劣最適性より
  \[
    \Wass_2(\mu_s, \mu_t)^2
    \leq \int_{\R^d} \norm{T_s(x) - T_t(x)}^2 \,\d\alpha(x).
  \]
  $T_s(x) - T_t(x) = (t - s)\bigl(x - \nabla\phi(x)\bigr)$ なので，$(\ast)$ より
  \[
    \Wass_2(\mu_s, \mu_t)^2
    \leq (t - s)^2 \int_{\R^d} \norm{x - \nabla\phi(x)}^2 \,\d\alpha(x)
    = (t - s)^2\, \Wass_2(\alpha, \beta)^2,
  \]
  すなわち $\Wass_2(\mu_s, \mu_t) \leq (t - s)\,\Wass_2(\alpha, \beta)$．

  \textbf{等号（測地線性）．}
  $\Wass_2$ の三角不等式（備考~\ref{rem:sem-w2-is-metric}）と上界より，
  $0 \leq s \leq t \leq 1$ に対し
  \[
    \Wass_2(\alpha, \beta)
    \leq \Wass_2(\mu_0, \mu_s) + \Wass_2(\mu_s, \mu_t) + \Wass_2(\mu_t, \mu_1)
    \leq \bigl(s + (t - s) + (1 - t)\bigr)\Wass_2(\alpha, \beta)
    = \Wass_2(\alpha, \beta).
  \]
  両端が等しいので途中の不等号はすべて等号でなければならず，
  とくに $\Wass_2(\mu_s, \mu_t) = (t - s)\,\Wass_2(\alpha, \beta)$．
  $s, t$ の大小によらず $|t-s|\,\Wass_2(\alpha,\beta)$ を得る．
\end{proof}

\begin{remark}{粒子の直線運動という描像}{sem-particle-picture}
  変位補間は，各点 $x$ の質量を直線
  $t \mapsto (1-t)x + t\,\nabla\phi(x)$ に沿って\textbf{一定速度}で
  動かす操作である．$L^2$ 補間
  $(1-t)\alpha + t\beta$（質量をその場で増減させる）とは異なり，
  $\Wass_2$ 測地線は質量を空間的に\textbf{輸送}する点が特徴である．
\end{remark}

\begin{remark}{変位凸性と曲率への布石}{sem-displacement-convexity-preview}
  測地線が定まると，汎関数 $F : \Pp_2(\R^d) \to \R$ の
  測地線に沿った凸性——\textbf{変位凸性}——を論じられる：
  $F$ が\textbf{変位凸}とは，任意の測地線 $(\mu_t)$ に沿って
  $t \mapsto F(\mu_t)$ が凸となることをいう．
  とくにエントロピー（自由エネルギー）汎関数の変位凸性は，
  $T = \nabla\phi$ のヤコビアン $\det(\partial^2\phi)$ の凹性
  （$\phi$ が凸なので $\partial^2\phi \succeq 0$）から導かれ，
  これが Ricci 曲率下界・CD$(K,N)$ 条件と結びつく．
  この計算では Brenier 写像が凸関数の勾配であること
  （主張~\ref{clm:sem-brenier}）が本質的に効く．
  動的定式化（Benamou--Brenier）と Otto 計算を経て，
  この曲率の主題へ進む．
\end{remark}
